In 1961, Claus Jönsson sent electrons through a double slit. In such an experiment, electrons accelerated through \UO\ pass a double slit with the spacing $d = \dsO$. The screen is $L = \LsO$ behind the slits; electron lenses then magnify the pattern.

\begin{abcliste}
\abc Calculate the de Broglie wavelength of the electrons.
\abc Calculate the spacing of neighbouring bright fringes on the screen.
\abc The voltage is made four times as large. How does the fringe spacing change?
\end{abcliste}
